\documentclass[11pt]{article}
\usepackage[a4paper,margin=27mm]{geometry}
\usepackage{amsmath,amssymb,mathtools}
\usepackage{fontspec,unicode-math}
\setmainfont{Latin Modern Roman}
\setmathfont{Latin Modern Math}
\usepackage{xurl,hyperref}
\hypersetup{colorlinks=true,urlcolor=blue,linkcolor=blue}
\newcommand{\sref}[2]{\hyperlink{source:#1:#2}{[S#2]}}
\newcommand{\cataloguescope}[1]{\paragraph{Mathematical question.} #1}
\begin{document}
% UnsolvedMath ID 8700020; source catalogue code AMR-086-0020
\setcounter{section}{8700019}
\section{Conjecture 3.2 — Roy}\label{8700020}
{\small\sffamily UnsolvedMath ID 8700020 · Graph Theory}
\subsection{Definitions and mathematical statement}

\paragraph{Shared notation.}$\mathbb N=\{1,2,\ldots\}$, and $\mathbb Z,\mathbb Q,\mathbb R,\mathbb C$ denote the integers, rationals, reals and complex numbers. Any other conventions are specified locally.


Fix a positive integer $r$, complex numbers $y_1,\ldots,y_r$ linearly independent over $\mathbb Q$, and nonzero complex numbers $\alpha_1,\ldots,\alpha_r$. On $\mathbb C[X_0,X_1]$ put $\mathcal D=\partial/\partial X_0+X_1\partial/\partial X_1$, and define $H(P)$ as the maximum absolute value of its coefficients. The fixed tuple size is denoted by $r$, while $j$ denotes a derivative order. Let $s_0,s_1,t_0,t_1,u>0$ satisfy \[\max\{1,t_0,2t_1\}<\min\{s_0,2s_1\},\qquad \max\{s_0,s_1+t_1\}<u<\tfrac12(1+t_0+t_1).\]
\paragraph{Statement recorded in the supplied source.}
Suppose that for every sufficiently large positive integer $N$ there is $0\ne P_N\in\mathbb Z[X_0,X_1]$ with $\deg_{X_0}P_N\le N^{t_0}$, $\deg_{X_1}P_N\le N^{t_1}$ and $H(P_N)\le e^N$, such that \[\left|(\mathcal D^jP_N)\left(\sum_{i=1}^r m_i y_i,\prod_{i=1}^r\alpha_i^{m_i}\right)\right|\le e^{-N^u}\] for every integer $0\le j\le N^{s_0}$ and every $m_1,\ldots,m_r\in\{0,1,2,\ldots\}$ with $\max_i m_i\le N^{s_1}$. Must \[\operatorname{trdeg}_{\mathbb Q}\mathbb Q(y_1,\ldots,y_r,\alpha_1,\ldots,\alpha_r)\ge r?\] \sref{8700020}{2}
\subsection{Short English statement}
Do uniformly small derivatives of bounded-degree integer polynomials along this additive–multiplicative grid force transcendence degree at least the fixed rank?
\subsection{Sources}
\begin{description}
\item[\hypertarget{source:8700020:1}{S1}] ULAM.AI and the UnsolvedMath Contributors, UnsolvedMath: A Curated Collection of Open Mathematics Problems, record 8700020 (AMR-086-0020). CC BY 4.0. \url{https://huggingface.co/datasets/ulamai/UnsolvedMath}.
\item[\hypertarget{source:8700020:2}{S2}] Michel Waldschmidt, Open Diophantine Problems, arXiv:math/0312440v2 (24 January 2004), Conjecture 3.2 and its preceding derivation and height definitions. \url{https://arxiv.org/abs/math/0312440}.
\end{description}
\label{end:8700020}
\end{document}
